\section{Lefschetz 纤维化的 Fukaya 范畴}\label{sec:18}\label{section:18}
\subsection{二重分歧覆盖}\label{sec:18a}
固定一个正合 Lefschetz 纤维化 $\pi:E^{2n+2}\to S$，其中 $S=D$ 为 $\CC$ 中带标准复结构的单位圆盘。选一个很小的 $\delta>0$，以 $S^{\mathrm{shrink}}\subset S$ 表示半径 $1-4\delta$ 的圆盘，以 $E^{\mathrm{shrink}}\subset E$ 表示它在总空间中的原像。假设 $\pi$ 的全部临界值都在 $S^{\mathrm{shrink}}$ 的内部；还假设，$E$ 上的负 Liouville 流经过某个正时间后，将整个 $E$ 推入 $E^{\mathrm{shrink}}$。当 $\delta$ 充分小时，这两个条件必定满足。取 $*=1-3\delta\in\operatorname{int}(S)$ 为基点，$M=E_*$ 为该点上的纤维。需要选择基路径 $\chi$ 时，始终使用从 $*$ 到 $1\in\partial S$ 的直线段。\znote{266}{四个半径的层级容易混淆，先列清楚：$S$ 是半径 $1$ 的圆盘；$\mathcal W=\{w:|w|\le1-2\delta\}$ 是 §18b 中族参数的取值圆盘；基点 $*=1-3\delta$ 落在 $\mathcal W$ 内但在 $S^{\mathrm{shrink}}$ 之外；而 $S^{\mathrm{shrink}}$ 与 $E^{\mathrm{shrink}}$ 的门槛是 $1-4\delta$。于是含 $*$ 的圆周 $|w|=1-3\delta$ 整条落在 $\mathcal W$ 内，引理 18.1 中的单值变换才是族 $\Psi$ 的合法环路；同时 $E^{\mathrm{shrink}}$ 上的一切估计都留出了 $\delta$ 的余量。后文还会出现第五个半径 $1-\delta$，用于定义 $\widetilde U$，不要与前四个混同。[来源：原书 §18a、§18b（印刷页 266、269）]}

设 $\widetilde S$ 为在 $*$ 处分歧的 $S$ 的二重覆盖；类似地，设 $\widetilde E$ 为沿 $M=E_*$ 分歧的 $E$ 的二重覆盖。当然，$\widetilde S$ 仍与 $D$ 同构，但我们
\SourcePage{274}
不作这一等同，而写成
\begin{equation}\label{eq:18.1}
\begin{aligned}
\widetilde S&=\{y\in\CC:y^2+*\in S\},\\
\widetilde E&=\{(y,x)\in\widetilde S\times E:y^2+*=\pi(x)\}.
\end{aligned}
\end{equation}
$\widetilde E$ 带有诱导的殆复结构，因为它是乘积 $\widetilde S\times E$ 的殆复子流形；这一结构使 $\partial\widetilde E$ 弱凸。关于该结构，有下列伪全纯映射：
\[
\begin{aligned}
\iota_{\widetilde S}(y)&=-y,&\iota_{\widetilde E}(y,x)&=(-y,x)&&\text{覆盖对合；}\\
\beta_{\widetilde S}(y)&=y^2+*:\widetilde S\to S,&\beta_{\widetilde E}(y,x)&=x:\widetilde E\to E&&\text{商映射；}\\
\widetilde\pi(y,x)&=y:\widetilde E\to\widetilde S&&&&\text{$\pi$ 的提升。}
\end{aligned}
\]
$\widetilde\pi$ 的临界点是 $\pi$ 的各临界点的原像；由于局部结构不变，它们仍然非退化。

为完成构造，还要引入适当的正合辛结构。回顾，定义纤维化的数据中，底空间 $S$ 带有某个 $\theta_S\in\Omega^1(S)$ 及 $\omega_S=d\theta_S\in\Omega^2(S)$。从\eref{18.1}可见，$\widetilde S$ 包含以原点为中心、半径 $\sqrt{2\delta}$ 的圆盘。设 $\kappa$ 为支撑在更小的半径 $\sqrt\delta$ 的圆盘内的一形式，关于 $\iota_{\widetilde S}$ 不变，且 $d\kappa$ 在原点处为正。我们将 $\kappa$ 看作很小的形式，因此保留在必要时将最初选取的形式乘以某个很小正常数的余地。令
\begin{equation}\label{eq:18.2}
\theta_{\widetilde S}=\beta_{\widetilde S}^*\theta_S+\kappa,\qquad
\omega_{\widetilde S}=d\theta_{\widetilde S}=\beta_{\widetilde S}^*\omega_S+d\kappa.
\end{equation}
由于 $d\kappa$ 很小，且在原点处为正，$\omega_{\widetilde S}$ 处处为正；又因为 $\kappa$ 支撑于内部，$\theta_{\widetilde S}$ 沿边界为正。这些条件保证\eref{18.2}使 $\widetilde S$ 成为带边界正合辛流形。与此相应，为 $\widetilde E$ 赋予
\begin{equation}\label{eq:18.3}
\theta_{\widetilde E}=\beta_{\widetilde E}^*\theta_E+\widetilde\pi^*\kappa,\qquad
\omega_{\widetilde E}=d\theta_{\widetilde E}=\beta_{\widetilde E}^*\omega_E+\widetilde\pi^*d\kappa.
\end{equation}
第一项 $\beta_{\widetilde E}^*\omega_E$ 只在分歧纤维 $\{y=0\}$ 的法向方向上不满足辛条件。因此，由与前面相同的论证，加入充分小的 $\widetilde\pi^*d\kappa$ 项，就使它处处为辛形式。注意，由于\eref{18.3}的第二项来自底空间，$\widetilde\pi$ 正则部分上的相应辛联络，是 $\pi$ 的辛联络的拉回。这说明 $\partial^h\widetilde E$ 是水平的；同时，使用\stref{引理}{15.1}中的相同论证，可知负 Liouville 流沿这一部分边界指向内部。在 $\partial^v\widetilde E$ 附近，负 Liouville 向量场是
\SourcePage{275}
$E$ 上相应向量场的拉回，因此也指向内部。这样，就验证了 $\widetilde\pi:\widetilde E\to\widetilde S$ 仍为正合 Lefschetz 纤维化。

在以下讨论中，$\widetilde E$ 中两类正合 Lagrangian 子流形将起突出作用。分别将它们记为 (U) 型与 (B) 型，对应于无分歧与有分歧：
\begin{definition*}[(U) 型]
覆盖 $\beta_{\widetilde S}$ 在 $S^{\mathrm{shrink}}$ 上的限制无分歧，因而平凡，也就是说，总空间由底空间的两个副本组成。$\beta_{\widetilde E}$ 在 $E^{\mathrm{shrink}}$ 上的限制也如此。记
\begin{equation}\label{eq:18.4}
\begin{aligned}
\widetilde S^{\mathrm{shrink}}&=\beta_{\widetilde S}^{-1}(S^{\mathrm{shrink}})=\widetilde S^{\mathrm{shrink}}_+\sqcup\widetilde S^{\mathrm{shrink}}_-,\\
\widetilde E^{\mathrm{shrink}}&=\beta_{\widetilde E}^{-1}(E^{\mathrm{shrink}})=\widetilde E^{\mathrm{shrink}}_+\sqcup\widetilde E^{\mathrm{shrink}}_-.
\end{aligned}
\end{equation}
更具体地说，满足 $|y^2+*|\le1-4\delta$ 的点 $y$ 必有 $\operatorname{im}(y)\ne0$，该虚部的符号区分了\eref{18.4}的两个分支。这些 $y$ 还满足 $|y|\ge\sqrt\delta$，因此在各二重分歧覆盖的限制版本上，\eref{18.2}、\eref{18.3}中的 $\kappa$ 项为零。\znote{267}{这两步都是作者的「显然」，补上推导：若 $y$ 为实数，则 $y^2+*\ge1-3\delta>1-4\delta$，与 $|y^2+*|\le1-4\delta$ 矛盾，故 $\operatorname{im}(y)\ne0$，\eref{18.4} 按虚部符号分成两支正由此而来；又 $|y^2|\ge|*|-|y^2+*|\ge(1-3\delta)-(1-4\delta)=\delta$，即 $|y|\ge\sqrt\delta$。因 $\kappa$ 支撑在 $\sqrt\delta$ 圆盘内，它在 $\widetilde S^{\mathrm{shrink}}$ 上恒为零，故\eref{18.2}、\eref{18.3}在该区域就是纯粹的拉回——这正是 (U) 型可写成两个不相交副本、且每支仍是正合 Lagrangian 的原因。[来源：原书 §18a（印刷页 267）]}一个显然的推论是：若 $L$ 为 $E$ 中包含于 $E^{\mathrm{shrink}}$ 的正合 Lagrangian 子流形，则它在 $\widetilde E$ 中的原像由 $L$ 的两个副本组成，且它们也都是正合 Lagrangian 子流形。类似于前面，记
\[
\widetilde L=\beta_{\widetilde E}^{-1}(L)=\widetilde L_+\sqcup\widetilde L_-.
\]
\end{definition*}
\begin{definition*}[(B) 型]
设 $\gamma:[0,1]\to S$ 为满足 $\gamma(0)=*$ 的消失路径，$\Delta_\gamma\subset E$ 为其 Lefschetz 顶针。那么
\begin{equation}\label{eq:18.5}
\widetilde\Delta_\gamma=\beta_{\widetilde E}^{-1}(\Delta_\gamma)
\end{equation}
为 $\widetilde E$ 中的光滑 Lagrangian 球面。可将其表示为匹配圈，从而看出这一点，具体如下。在相差定向反转的意义下，存在唯一的光滑路径 $\widetilde\gamma:[-1,1]\to\widetilde S$，满足 $\beta_{\widetilde S}(\widetilde\gamma(t))=\gamma(t^2)$。$\widetilde\gamma$ 的端点是临界值 $\gamma(1)$ 的原像，因此其两个半段 $\widetilde\gamma_\pm(t)=\widetilde\gamma(\pm t):[0,1]\to\widetilde S$ 均为 $\widetilde\pi$ 的消失路径。如上所述，辛联络是 $\pi$ 的辛联络的拉回，由此可知，$\widetilde\gamma_\pm$ 的每个 Lefschetz 顶针都投影到 $\gamma$ 的 Lefschetz 顶针。这意味着，$\widetilde\gamma_\pm$ 的消失圈重合，等于 $M$ 内 $\gamma$ 的消失圈；而\eref{18.5}就是 $\widetilde\gamma$ 的朴素意义下的匹配圈，见\secref{16g}开头。
\end{definition*}
可以用另一个正则值 $w\in S\setminus\operatorname{Critv}(\pi)$ 代替 $*=1-3\delta$ 作为分歧点。证明二重分歧覆盖为 Lefschetz 纤维化的技术细节要求 $|w|\le1-2\delta$。构造 (U) 型 Lagrangian 子流形时，须假设 $L$ 落在
\SourcePage{276}
$|\pi(x)|\le|w|-\delta$ 的区域内。此外，$\widetilde L_+$ 与 $\widetilde L_-$ 的区分依赖于 $\sqrt w$ 的选择。当然，这意味着，若令 $w$ 沿例如半径 $1-3\delta$ 的圆周移动，这两个分支就会互换。下面将对此作严格说明，见\stref{引理}{18.1}。
